%%%%%%%%%%%%%%%%%%%%%%%%%%%%%%%%%%%%%%%%%%%%%%%%%%%%%%%%%%%%%%%%%%%%
%% I, the copyright holder of this work, release this work into the
%% public domain. This applies worldwide. In some countries this may
%% not be legally possible; if so: I grant anyone the right to use
%% this work for any purpose, without any conditions, unless such
%% conditions are required by law.
%%%%%%%%%%%%%%%%%%%%%%%%%%%%%%%%%%%%%%%%%%%%%%%%%%%%%%%%%%%%%%%%%%%%

\documentclass{beamer}
\usetheme[faculty=ped,basePath=../fibeamer,logo=../fibeamer/logo/mu/Logo-UC-3.png]{fibeamer}
\graphicspath{{../resources/}}

\usepackage{algorithm}
\usepackage{algpseudocode}

\usepackage[utf8]{inputenc}
\usepackage[
  main=english, %% By using `czech` or `slovak` as the main locale
                %% instead of `english`, you can typeset the
                %% presentation in either Czech or Slovak,
                %% respectively.
  czech, slovak %% The additional keys allow foreign texts to be
]{babel}        %% typeset as follows:
%%
%%   \begin{otherlanguage}{czech}   ... \end{otherlanguage}
%%   \begin{otherlanguage}{slovak}  ... \end{otherlanguage}
%%
%% These macros specify information about the presentation
\title{Funciones propias} %% that will be typeset on the
\subtitle{Semestre II-2026} %% title page.
\author{Profesora: Andreina Cota Vidaurre}
%% These additional packages are used within the document:
\usepackage{ragged2e}  % `\justifying` text
\usepackage{booktabs}  % Tables
\usepackage{tabularx}
\usepackage{tikz}      % Diagrams
\usetikzlibrary{calc, shapes, backgrounds}
\usetikzlibrary{positioning, arrows.meta, shadows.blur, shapes.geometric}
\usepackage{amsmath, amssymb}
\usepackage{url}       % `\url`s
\usepackage{listings}  % Code listings
\usepackage{graphicx}
\usepackage{amsmath}
\usepackage[T1]{fontenc}
\usepackage{xcolor}
\usepackage{minted}
\usepackage{pifont}
% \usepackage{adjustbox}


% Definicion de pseudocodigo en espaniol
\lstdefinelanguage{PseudocodigoES}{
    morekeywords={INICIO,FIN,PARA,HACER,SI,ENTONCES,FIN_SI,FIN_PARA,IMPRIMIR},
    sensitive=true,
    morecomment=[l]{//},
    morestring=[b]"
}

% Definicion de pseudocodigo en ingles
\lstdefinelanguage{PseudocodeEN}{
    morekeywords={BEGIN,END,FOR,DO,IF,THEN,END_IF,END_FOR,PRINT},
    sensitive=true,
    morecomment=[l]{//},
    morestring=[b]"
}

\lstset{
    basicstyle=\ttfamily\small,
    keywordstyle=\bfseries,
    columns=fullflexible,
    keepspaces=true,
    showstringspaces=false,
    frame=single,
    xleftmargin=0em,
    framexleftmargin=0em,
    framesep=2pt,
    gobble=4,
    resetmargins=true,
    aboveskip=0pt,
    belowskip=0pt
}

\lstdefinestyle{pythonstyle}{
    language=Python,
    basicstyle=\ttfamily\small,
    keywordstyle=\bfseries,
    commentstyle=\itshape,
    stringstyle=\ttfamily,
    showstringspaces=false,
    columns=fullflexible,
    keepspaces=true,
    frame=single,
    breaklines=true,
    xleftmargin=0em,
    framexleftmargin=0em,
    framesep=2pt,
    aboveskip=0pt,
    belowskip=0pt,
    gobble=4,
    resetmargins=true,
    emph={print,input,int,float,str,bool,len,range},
    emphstyle=\bfseries
}

\lstset{style=pythonstyle}

% Estilos del diagrama
\tikzstyle{iniciofin} = [
    ellipse,
    minimum width=2.5cm,
    minimum height=1cm,
    text centered,
    draw=black,
    fill=blue!25
]

\tikzstyle{proceso} = [
    rectangle,
    rounded corners,
    minimum width=3cm,
    minimum height=1cm,
    text centered,
    draw=black,
    fill=cyan!20
]

\tikzstyle{decision} = [
    diamond,
    aspect=2,
    text centered,
    draw=black,
    fill=yellow!35
]

\tikzstyle{flecha} = [
    thick,
    ->,
    >=Stealth
]

% \definecolor{green_python}{HTML}{52A736}

% \newcolumntype{Y}{>{\raggedright\arraybackslash}X}

\frenchspacing
\begin{document}
  \shorthandoff{-}
  \frame[c]{\maketitle}

% \begin{darkframes}

    \begin{frame}[fragile, t]{Contenido}
        \begin{itemize}
            \item ¿Qué es una función?
            \item Función vs método
            \item Estructura y sintaxis de una función
            \item Parámetros y argumentos
            \item Tipos de funciones
            \item Ejemplos
            \item Resumen y aspectos importantes
        \end{itemize}
    \end{frame}

    \begin{frame}[fragile, t]{Qué es una función}
        \begin{itemize}
            \item Una función es un bloque de código reutilizable que realiza una tarea específica. 
            \item Se define una vez y se puede ejecutar (llamar) cuantas veces sea necesario.
            \item Tiene un nombre que la identifica
            \item Puede recibir información de entrada
            \item Puede devolver un resultado
            \item Evita repetir código
        \end{itemize}
    \end{frame}

    \begin{frame}{Qué es una función}
        \begin{itemize}
            \item Una función es como una receta de cocina: se escribe una sola vez, pero cualquier persona puede seguirla en cualquier momento y obtener el mismo resultado.
            \item Una función define claramente qué información se necesita (parámetros), qué pasos ejecutar (cuerpo) y qué resultado reportar (retorno).
        \end{itemize}
    \end{frame}

    \begin{frame}[fragile,t]{Funciones vs métodos}
        \begin{itemize}
            \item Función:
            \begin{itemize}
                \item Es independiente; no pertenece a ningún objeto
                \item Se llama directamente por su nombre
                \item Ejemplo: \mintinline{python}{len("hola"), print("UC"), max(3, 7, 1)}
            \end{itemize}
            \item Método:
                \begin{itemize}
                    \item Es una función que pertenece a un objeto o clase
                    \item Se llama a través del objeto
                    \item Ejemplo: \mintinline{python}{"UC".find("C"), "  carrera  ".strip()}
                \end{itemize}
        \end{itemize}
        \textbf{Regla simple:} Todo método es una función, pero no toda función es un método.
    \end{frame}

    \begin{frame}[fragile, t]{Función vs método}
        \begin{tabular}{|p{2.5cm}|p{3.3cm}|p{3.3cm}|}
            \hline
             & Función & Método \\ \hline
            ¿A quién pertenece? & A nadie / al programa & A un objeto o clase \\ \hline
            ¿Cómo se llama? & \mintinline{python}{nombre()} & \mintinline{python}{objeto.nombre()} \\ \hline
            Contexto & Programación general & Progranación orientada a objetos \\
            \hline
        \end{tabular}
    \end{frame}

    \begin{frame}[fragile,t]{Estructura y sintaxis de una función}
        \begin{minted}[xleftmargin=-2cm, frame=leftline, fontsize=\small]{python}
            def nombre_funcion(parametro1, parametro2):
                """Descripción de lo que hace la función"""
                # Cuerpo de la función
                resultado = parametro1 + parametro2
                return resultado
        \end{minted}
    \end{frame}

    \begin{frame}[fragile, t]{Estructura y sintaxis de una función}
        \centering
        \begin{tabular}{|p{4.5cm}|p{5cm}|}
            \hline
            \mintinline{python}{def} & Palabra reservada que indica que se define una función \\ \hline
            \mintinline{python}{nombre_funcion()} & Identificador único de la función \\ \hline
            \mintinline{python}{(parametro1, parametro2)} & Datos de entrada que puede recibir \\ \hline
            \mintinline{python}{"""..."""} & Documentación opcional pero recomendada \\ \hline
            Cuerpo & Instrucciones que ejecuta la función \\ \hline
            \mintinline{python}{return} & Valor que devuelve al terminar (opcional) \\ \hline
        \end{tabular}
    \end{frame}
    
    \begin{frame}[fragile, t]{Parámetros y argumentos}
        \begin{itemize}
            \item \textbf{Parámetro}: Variable que se declara en la definición de la función. Es un marcador de posición.
            \item \textbf{Argumento}: Valor real que se pasa a la función al llamarla.
        \end{itemize}
        \begin{minted}[xleftmargin=-2cm, frame=leftline, fontsize=\small]{python}   
            # "nombre" y "seccion" son PARÁMETROS
            def presentar_estudiante(nombre, seccion):
                print(f"El estudiante {nombre} esta en la seccion {seccion}.")
            
            # "Juan" y "12" son ARGUMENTOS
            presentar_estudiante("Juan", 12)
        \end{minted}
    \end{frame}

    \begin{frame}[fragile, t]{Tipos de parámetros}
        \begin{itemize}
            \item Requeridos: deben entregarse siempre
            \begin{minted}[xleftmargin=-2cm, frame=leftline, fontsize=\small]{python}
                "PUC".find("U")
            \end{minted}
            \item Con valor por defecto: tienen un valor si no se entrega ninguno
            \begin{minted}[xleftmargin=-2cm, frame=leftline, fontsize=\small]{python}
                range(1, 11, 2) # start, stop, step
            \end{minted}
            \item Arbitrarios \mintinline{python}{*args}: aceptan cantidad variable de argumentos
            \begin{minted}[xleftmargin=-2cm, frame=leftline, fontsize=\small]{python}
                print("UC", "Facultad")
            \end{minted}
        \end{itemize}
    \end{frame}
    
    \begin{frame}[fragile, t]{Tipos de funciones propias}
        \begin{itemize}
            \item Sin parámetros y sin retorno
            \begin{minted}[xleftmargin=-2cm, frame=leftline, fontsize=\small]{python}
            def dar_saludo():
                print("Bienvenido a la UC")
            \end{minted}
            \item Con parámetros y sin retorno
            \begin{minted}[xleftmargin=-2cm, frame=leftline, fontsize=\small]{python}
            def mostrar_nombre(nombre):
                print("Nombre: " + nombre)
            \end{minted}
            \item Sin parámetros y con retorno
            \begin{minted}[xleftmargin=-2cm, frame=leftline, fontsize=\small]{python}
            def obtener_nivel_alerta():
                return "AMARILLO"
            \end{minted}
            \item Con parámetros y con retorno (la más completa)
            \begin{minted}[xleftmargin=-2cm, frame=leftline, fontsize=\small]{python}
            def calcular_nota(nota1, nota2):
                total = nota1 + nota2
                return total
            \end{minted}
        \end{itemize}
    \end{frame}

    \begin{frame}[fragile, t]{Ejemplos de funciones}
        Preparar café
        \begin{minted}[xleftmargin=-2cm, frame=leftline, fontsize=\small]{python}
            def preparar_cafe(tipo, azucar=1):
                print(f"Preparando café {tipo} con {azucar} cucharada(s)")
            
            preparar_cafe("americano")
            preparar_cafe("espresso", 0)
        \end{minted}
    \end{frame}

    \begin{frame}[fragile, t]{Ejemplos de funciones}
        Calcular propina
        \begin{minted}[xleftmargin=-2cm, frame=leftline, fontsize=\small]{python}
            def calcular_propina(total, porcentaje=10):
                propina = total * (porcentaje / 100)
                return round(propina, 2)
            
            print(calcular_propina(25000))       # Propina del 10%
            print(calcular_propina(25000, 15))   # Propina del 15%
        \end{minted}
    \end{frame}

    \begin{frame}[fragile, t]{Ejemplos de funciones}
        Enviar un mensaje
        \begin{minted}[xleftmargin=-2cm, frame=leftline, fontsize=\small]{python}
            def enviar_mensaje(destinatario, texto):
                print("Enviando a " + destinatario + ": '" + texto + "'")
            
            enviar_mensaje("Pedro", "El examen es a las 18:00")
        \end{minted}
    \end{frame}

    % \begin{frame}[fragile, t]{Ejemplos de funciones en el ámbito militar}
    %     Reporte de situación (SITREP)
    %     \begin{minted}[xleftmargin=-2cm, frame=leftline, fontsize=\small]{python}
    %         def emitir_sitrep(unidad, posicion, estado):
    %             print(f"Unidad:{unidad}|Pos:{posicion}|Estado:{estado}")
            
    %         emitir_sitrep("Alpha-1", "Grid 447823", "OPERATIVA")
    %         emitir_sitrep("Bravo-3", "Grid 112045", "CONTACTO")
    %     \end{minted}
    % \end{frame}

    % \begin{frame}[fragile, t]{Ejemplos de funciones en el ámbito militar}
    %     Calcular autonomía de combustible
    %     \begin{minted}[xleftmargin=-2cm, frame=leftline, fontsize=\small]{python}
    %         def autonomia(litros_disponibles, consumo_por_km):
    %             km_posibles = litros_disponibles / consumo_por_km
    %             return round(km_posibles, 1)
            
    %         print(autonomia(200, 0.4))  # Retorna 500.0 km
    %     \end{minted}
    % \end{frame}

    % \begin{frame}[fragile, t]{Ejemplos de funciones en el ámbito militar}
    %     Verificar acceso a instalación
    %     \begin{minted}[xleftmargin=-2cm, frame=leftline, fontsize=\small]{python}
    %         def verificar_acceso(id, nivel_requerido, nivel_personal):
    %             if nivel_personal >= nivel_requerido:
    %                 print(f"Acceso AUTORIZADO para ID {id}")
    %             else:
    %                 print(f"Acceso DENEGADO para ID {id}")
            
    %         verificar_acceso("MIL-0042", 3, 4)  # Autorizado
    %         verificar_acceso("MIL-0091", 3, 2)  # Denegado
    %     \end{minted}
    % \end{frame}

    \begin{frame}[fragile,t]{Buenas prácticas}
        \begin{enumerate}
            \item Nombres descriptivos
            \begin{minted}[xleftmargin=-2cm, frame=leftline, fontsize=\small, escapeinside=||]{python}
            # |\textcolor{red}{\ding{55}}| Malo
            def f(x, y):
                return x + y
            
            # |\textcolor{green!60!black}{\ding{51}}| Bueno
            def calcular_total_nota(nota1, nota2):
                return nota1 + nota2
            \end{minted}
            \item Una función = una responsabilidad: Cada función debe hacer UNA sola cosa y hacerla bien.
        \end{enumerate}
    \end{frame}

    \begin{frame}[fragile,t]{Buenas prácticas}
        \begin{enumerate}[3]
            \item Documentar con docstrings
            \begin{minted}[xleftmargin=-2cm, frame=leftline, fontsize=\small]{python}
            def calcular_autonomia(litros, consumo):
                """
                Calcula los kilómetros posibles según combustible 
                disponible.
                """
                return litros / consumo
            \end{minted}
        \end{enumerate}
    \end{frame}

    \begin{frame}[fragile, t]{Resumen y aspectos importantes}
        \centering
        \begin{tabular}{|p{3cm}|p{6cm}|}
            \hline
            Concepto & En una línea \\ \hline
            Función & Bloque de código reutilizable con nombre propio \\ \hline
            Método & Función que pertenece a un objeto o clase \\ \hline
            Parámetro & Función que pertenece a un objeto o clase \\ \hline
            Argumento & Valor real entregado al llamar la función \\ \hline
            \mintinline{python}{return} & Permite que la función devuelva un resultado \\ \hline
            \mintinline{python}{def} & Palabra reservada para definir una función en Python \\ \hline
        \end{tabular}
    \end{frame}

    \begin{frame}{Razones para usar funciones}
        \begin{itemize}
            \item \textbf{Reutilización}: Escribe una vez, usa muchas veces
            \item \textbf{Organización}: Código más limpio y ordenado
            \item \textbf{Mantenimiento}: Si hay un error, se corrige en un solo lugar
        \end{itemize}
        
        \textbf{"Una buena función hace una cosa, la hace bien y tiene un nombre que lo dice."}
    \end{frame}

\end{document}
